\section{Method}

\subsection{Planned Method}

The original plan was to train a DT on optimizer-generated trajectories and then use autoregressive rollout to initialize the nonlinear solver. The hoped-for outcome was a warm-start that was already close to a feasible raceline around obstacles, reducing solver iterations and total solve time. Post-projection repair data and later failure-conditioned data were proposed specifically to address rollout distribution shift: once the DT leaves the expert manifold, optimizer-labeled corrections may provide the right training signal.

\subsection{Actual Implemented Method}

The implemented DT uses a causal transformer over triplets of return-to-go, augmented state, and action. The augmented state contains the vehicle/path observation
\[
[u_x, u_y, r, e, d\psi, \mathrm{pos}_E, \mathrm{pos}_N, \psi_{\mathrm{world}}],
\]
track features $(\kappa, \text{half-width})$, and a fixed number of obstacle slots containing lookahead obstacle features. The main model predicts the current action, and training also uses an auxiliary next-state prediction head.

Training minimizes
\[
\mathcal{L} = \mathcal{L}_{a} + \lambda_x \mathcal{L}_{x},
\]
where both losses are masked MSE terms normalized only over valid, non-padding tokens. This detail matters because the repo explicitly fixed masked-loss normalization and padded-window handling during the project. Later fine-tunes also used state-dimension weighting that emphasized $u_y$, $r$, $e$, and $d\psi$.

At inference time, the warm-starter rolls out the DT autoregressively. Each step:
\begin{enumerate}
\item builds the augmented observation from the current state, track geometry, and nearby obstacles;
\item predicts the next control action;
\item propagates dynamics to get the next state;
\item optionally projects the next state back toward a safe envelope; and
\item validates the resulting rollout before handing it to the optimizer.
\end{enumerate}

The wrapper supports three projection modes:
\texttt{off}, \texttt{soft}, and \texttt{full}. It also tracks fallback counts, projection fraction, total projection magnitude, max projection magnitude, and per-reason counts such as \texttt{e\_clip} and \texttt{obs\_push}. These diagnostics became central to interpreting downstream failures.

\subsection{Wrapper Behavior and Acceptance Gate}

The repo's deployment logic does not automatically trust any DT rollout. A generated warm-start is accepted only if it passes both geometric validation and a threshold-based gate. In the default diagnostic configuration, the gate rejects a rollout if any of the following hold:
\begin{itemize}
\item fallback count $> 0$,
\item projection fraction $> 0.8$,
\item projection total magnitude $> 100$,
\item projection max-step magnitude $> 2.0$.
\end{itemize}

This gate is important for two reasons. First, it protects the optimizer from poor initializations. Second, it means that offline improvement alone is insufficient: the DT must produce rollouts that are good enough under the actual acceptance policy.
